\begin{result}{F.3.24}{The genus-three universal quotient on balanced tangencies}{Proved; geometric monodromy and exact trace equations}
Let $V$ be a symplectic vector space of dimension six over
$\mathbf Q_\ell$, with ordered symplectic basis
$a,b,c,d,e,f$. Put $z_1=[a,b]$, $z_2=[c,d]$,
$z_3=[e,f]$, and $\omega=z_1+z_2+z_3$ in the free Lie
algebra $\mathbb L(V)$. Let $\mathfrak h_6$ be its special
degree-six derivations. Give the basis the characters
\[
 a:\alpha_1,\ b:t/\alpha_1,\quad
 c:\alpha_2,\ d:t/\alpha_2,\quad
 e:\alpha_3,\ f:t/\alpha_3
\]
of the full diagonal symplectic-similitude torus. Define
$T\subset\mathfrak h_6$ to be its $t^3$-weight subspace
of derivations tangent to every one of the six basis lines:
\[
 D(v)\in[\mathbb L_6(V),v]\qquad(v=a,b,c,d,e,f).
\]
Let $U_6\subset\mathfrak h_6$ be the rational universal
arithmetic degree-six monodromy image, with the unique
special representative used in F.3.22, and let
$q_6:U_6\twoheadrightarrow Q_6$ be Nakamura's universal
arithmetic quotient, whose similitude character is $t^3$\cite{AuditNakamura1996}.
Then
\[
                         q_6(T\cap U_6)=0.                 \tag{G3.1}
\]
This is a statement about the specified universal maps.
It does not assert that an arbitrary detecting cover in
LM6 A.22 supplies the full diagonal torus or all six
simultaneous basis-line constraints.

\end{result}
\begin{proof}
We identify the completed geometric image before imposing the finite trace equations.
\paragraph{1. The exact comparison with geometric monodromy.}
Use the moduli stack of genus-three curves equipped with
one marked point and a nonzero tangent vector there.
Its geometric fundamental group is the profinite
completion of the mapping class group of a genus-three
surface with one boundary circle fixed pointwise.
The universal tangential basepoint gives a monodromy
homomorphism into automorphisms of the punctured
fundamental group fixing the boundary word geometrically.
The arithmetic action sends this word to its cyclotomic
power. These are actual homomorphisms from the exact
arithmetic fundamental-group sequence of this moduli stack.
In particular its geometric monodromy image is normal
in the full arithmetic image.

The tangential arithmetic action can also be described
on each number-field fibre without choosing a rational
point of the open curve. A local uniformizer $t$ whose
first derivative agrees with the chosen tangent defines
the coefficient action of $G_K$ on
$\overline K((t^{1/n}))$ for all compatible $n$th roots,
fixing $t^{1/n}$. The resulting action on tame inertia
is the cyclotomic action. Replacing the coordinate by
$t+O(t^2)$ gives the compatible root substitutions with
leading coefficient one, which identify these tangential
fibre functors. This is the marked-tangent version of
the universal action.

Write $R$ for the proalgebraic Zariski closure over
$\mathbf Q$ of the discrete geometric monodromy in
automorphisms of the rational Malcev Lie algebra of
the punctured surface. At every finite nilpotent level,
the actual profinite geometric image lies in
$R(\mathbf Q_\ell)$: its dense discrete subgroup does,
and all defining polynomials are continuous over
$\mathbf Q_\ell$. Zariski closure commutes with this
scalar extension. Indeed, in every fixed degree the
linear equations imposed on the coefficients of a
vanishing polynomial have rational entries; their
solution space after scalar extension is exactly the
extended rational solution space.

The precise identification of the graded Lie algebra
of this image comes from the following opened primary
locators in Hain's final 1997 article\cite{AuditHain1997Torelli}. Theorem~3.4,
printed p.~605, makes the prounipotent radical of
relative completion a quotient of the Malcev
completion of the Torelli group. Lemma~4.5 and
Corollary~4.8, pp.~610--611, give the compatible
action and its weight filtration. The punctured
surface and boundary-fixed action are treated
explicitly in Section~14.4, p.~643: all the action
and filling maps are morphisms of mixed Hodge
structures, hence strict. The image of the relative
completion is $R$ by its universal property and
Zariski density. More explicitly, Theorem~4.10,
p.~611, identifies the Torelli Malcev weight
filtration with its lower central series, and
Remark~4.11, p.~613, records the compatible action
on the Malcev Lie algebras. Strictness of the
punctured action therefore makes its graded image
generated by the degree-one Johnson image.
Enomoto--Satoh, Theorem~3.5 and Proposition~3.6,
p.~640, identify that generated algebra with
the actual rational Johnson image. Consequently
\[
 \operatorname{gr}_{-m}^W\operatorname{Lie}R_u
   =\operatorname{Im}\bigl(\tau_{g,1}(m)\otimes\mathbf Q\bigr)
   =:J_m,                                                  \tag{G3.2}
\]
as subspaces of the actual derivation space. Here the
weight filtration on the image is the filtration
induced from its action on the punctured Malcev Lie
algebra. Strictness is essential in this assertion.
This argument does not identify closures of individual
discrete filtration kernels by density alone.

The canonical special representative does not destroy
normalization of $R$. For an arithmetic element trivial
on homology and in outer filtration degree six, its
tangential action fixes the full boundary logarithm
$Z$, whose initial term is $\omega$. Any lower-degree
inner part has degree-two term $c\omega$ and no
degree-one term. Remove it by $\exp(-c\operatorname{ad}Z)$.
This is a point of the geometric algebraic image: the
boundary Dehn twist has logarithm
$\operatorname{ad}Z$, and its integer powers are
Zariski dense in the corresponding additive one-parameter
group. It therefore preserves normalization of $R$.
The free-Lie centralizer of $\omega$ then removes
degrees three through five exactly as in F.3.22.
The resulting degree-six derivation is the unique
special representative. It follows, by taking the
leading term of its action on $\operatorname{Lie}R$, that
\[
                       [D,J_1]\subset J_7
                       \qquad(D\in U_6).                  \tag{G3.3}
\]
The commutator used here is an actual normalizer action
on the algebraic image, followed by its strict graded
map (G3.2).

\paragraph{2. The trace equation and its precise genus range.}
The degree-one Johnson image is
$J_1=\Lambda^3V$, with the derivation corresponding to
$u\wedge v\wedge w$ given by
\[
 \delta_{uvw}(x)=\langle x,u\rangle[v,w]
                +\langle x,v\rangle[w,u]
                +\langle x,w\rangle[u,v].                 \tag{G3.4}
\]
Define $\operatorname{ES}_k$ on a degree-$k$ special
derivation by the exact composition
\[
 \mathfrak h_k\hookrightarrow V\otimes\mathbb L_{k+1}(V)
 \hookrightarrow V^{\otimes(k+2)}
 \xrightarrow{\,\text{contract slots }1,2\,}V^{\otimes k}
 \longrightarrow (V^{\otimes k})_{\mathbf Z/k}.             \tag{G3.5}
\]
The tensor representing $D$ is
$\sum_i(x_i\otimes D(y_i)-y_i\otimes D(x_i))$.
Changing the simultaneous duality convention changes
only the overall sign and does not change the kernel.

Enomoto--Satoh's final 2014 article\cite{AuditEnomotoSatoh2014}, Theorem~3.5 and
Proposition~3.6 on p.~640 and the actual diagram in
Section~7.1, p.~656, give
$\operatorname{ES}_7(J_7)=0$ in their stable genus
range. To apply it here, attach six additional handles,
giving the actual inclusion of boundary-fixed mapping
class groups $\Gamma_{3,1}\to\Gamma_{9,1}$ and the
symplectic inclusion $V\hookrightarrow V'$. Extend a
derivation by zero on the added generators. Both the
Johnson map and (G3.5) commute with this inclusion.
The map on cyclic tensor coinvariants is injective:
projection killing the added letters is its left
inverse. Since $9\geq7+2$, the stated stable theorem
therefore gives the required vanishing at genus three.
Combining it with (G3.3) gives
\[
       \operatorname{ES}_7([D,\delta])=0
       \quad(D\in U_6, \delta\in\Lambda^3V).              \tag{G3.6}
\]
No stable-range formula is directly substituted at
genus three without these explicit maps.

\paragraph{3. The complete tangent space.}
Every element of $T$ has a unique expression
$D(v)=[A_v,v]$ with $A_v$ of degree six and character
$t^3$. Uniqueness follows from the centralizer of a
free generator, and the character assertion follows
by the injectivity of $A_v\mapsto[A_v,v]$ in this degree.
Its balanced multidegrees have three pairs and are
indexed by $(m_1,m_2,m_3)$ with sum three.
The specialness equation $D\omega=0$ splits into
blocks indexed by $(n_1,n_2,n_3)$ with sum four.
Up to permutation they have the following sizes and
rational ranks:
\[
\begin{array}{c|c|c|c}
 (n_1,n_2,n_3)&\text{rows}&\text{columns}&\text{rank}\\\hline
 (4,0,0)&8&6&6\\
 (3,1,0)&140&66&66\\
 (2,2,0)&312&120&119\\
 (2,1,1)&1260&360&359.
\end{array}                                               \tag{G3.7}
\]
The first three blocks are exactly the two-handle
blocks certified in F.3.22, with absent letters omitted.
For the new block a nonzero $359$-minor modulo $101$
has determinant $100$, and the following exact integer
kernel provides the matching upper bound.

For each pair $i<j$, put $C_{ij}=[z_i,z_j]$ and define
$P_{ij}$ by common conjugators on each handle:
\[
 A_i=[z_j,C_{ij}],\qquad A_j=-[z_i,C_{ij}],\qquad A_k=0.
\]
For each $i$, with $\{i,j,k\}=\{1,2,3\}$, define $T_i$ by
\[
 \begin{split}
 A_i&=[z_j,[z_i,z_k]]+[z_k,[z_i,z_j]],\\
 A_j&=-[z_i,[z_i,z_k]],\qquad
 A_k=-[z_i,[z_i,z_j]].
 \end{split}                                               \tag{G3.8}
\]
Each handle's two generators have the same displayed
conjugator. Jacobi gives $\sum[A_i,z_i]=0$ exactly.
The six elements
\[
                  P_{12},P_{13},P_{23},T_1,T_2,T_3
\]
are therefore a basis of $T$ by (G3.7).

\paragraph{4. The exact normalizer trace kernel.}
Apply (G3.5) to brackets with the twenty derivations
(G3.4) indexed by triples of distinct basis letters.
Using cyclic words of length seven as coordinates
gives an integer matrix with $2208$ nonzero rows and
six columns in the order above. Four of its rows,
labelled by $(\delta,\text{cyclic word})$, are
\[
 (abc,abccdcd),\quad(abc,aabcfeb),\quad
 (abc,abccdef),\quad(abe,aabcdeb).
\]
On the first four columns their matrix is
\[
 \begin{pmatrix}
 -1&0&0&0\\
 0&-1&0&1\\
 0&0&1&-1\\
 -1&0&0&1
 \end{pmatrix},\qquad\det=1.                              \tag{G3.9}
\]
Exact expansion of all the bracket traces gives identical
columns for $T_1,T_2,T_3$. Thus the full kernel, with
both bounds certified over the integers, is
\[
 \{D\in T:\operatorname{ES}_7([D,\Lambda^3V])=0\}
       =\mathbf Q_\ell(T_2-T_1)+\mathbf Q_\ell(T_3-T_1).
                                                               \tag{G3.10}
\]
The complete code keeps the exact associative-word
expansions; the minor alone supplies only the lower
bound and is used together with the full-column equality.

\paragraph{5. The arithmetic quotient.}
Permuting the three symplectic handles conjugates the
three derivations $T_i$ transitively. Every
$\operatorname{Sp}(V)$-invariant linear functional on
$\mathfrak h_6$ therefore vanishes on both differences
in (G3.10). The universal image $U_6$ is an algebraic
$\operatorname{Sp}(V)$-submodule: conjugation by the
geometric mapping-class image preserves it, and that
image is Zariski dense in $\operatorname{Sp}(V)$ on
homology. Nakamura's equivariance makes $q_6$ invariant
under this symplectic group, since its target character
is $t^3$. Complete reducibility in characteristic zero
extends this invariant linear map from $U_6$ to
$\mathfrak h_6$ (or, equivalently, gives an equivariant
projection $\mathfrak h_6\to U_6$). Equations (G3.6)
and (G3.10) now imply (G3.1).

\paragraph{Scope and the failure of the simpler filling calculation.}
The two-handle derivation $P_{12}$ has nonzero image on
the proper genus-three quotient. The actual free-Lie
substitution
\[
 a,c,e\mapsto x,\qquad b\mapsto y,\qquad d\mapsto z,
 \qquad f\mapsto-y-z
\]
kills $\omega$, whereas the image of $P_{12}(a)$ has
coefficient one on the word $xxyxzxz$. This proves the
nonvanishing through an explicit quotient morphism.
The additional geometric-normalizer equation (G3.6)
is what removes the relevant arithmetic quotient.
No element of a common-kernel subgroup with an open
arithmetic normalizer is constructed or ruled out on
an arbitrary LM6 A.22 detecting cover by this lemma
alone. Its exact further requirements are the actual
CM-character projection and enough simultaneous
orbit-line constraints on that cover.

\end{proof}

